% Acronyms
%\newacronym{api}{API}{Application Program Interface}
% \newacronym{tsa}{TSA}{Termine solo acronimo}

\newacronym{erp}{ERP}{Enterprise Resource Planning}
\newacronym{mvvm}{MVVM}{Model-View-ViewModel}
\newacronym{uml}{UML}{Unified Modeling Language}
\newacronym{ia}{IA}{Intelligenza Artificiale}


% Glossary
% \newglossaryentry{TermineSenzaAcronimo}{
%     name={Nome del termine},
%     sort=termine senza acronimo,
%     description={Descrizione}
% }

\newglossaryentry{apig}{
    name={API},
    text={Application Program Interface},
    sort=api,
    description={In informatica, un'API  è un insieme di regole e protocolli che consente a diverse applicazioni software di comunicare tra loro. Le API definiscono come le richieste e le risposte devono essere formattate, quali dati possono essere scambiati e quali operazioni possono essere eseguite. Le API sono fondamentali per l'integrazione di sistemi e lo sviluppo di applicazioni moderne, poiché permettono agli sviluppatori di utilizzare funzionalità preesistenti senza doverle implementare da zero.}
}

\newglossaryentry{serverconsolidationg}{
    name={Server Consolidation},
    text={Server Consolidation},
    sort=server consolidation,
    description={Server Consolidation è il processo di riduzione del numero di server fisici in un'organizzazione, al fine di ottimizzare l'utilizzo delle risorse e migliorare la gestione dell'infrastruttura.}
}

\newglossaryentry{erpg}{
    name={ERP},
    text={Enterprise Resource Planning},
    sort=erp,
    description={ERP è un sistema software integrato che gestisce e automatizza i processi aziendali, come contabilità, gestione delle risorse umane, produzione, logistica e vendite. L'obiettivo principale di un ERP è migliorare l'efficienza operativa e la visibilità dei dati all'interno dell'organizzazione, consentendo una migliore pianificazione e controllo delle risorse.}
}

\newglossaryentry{ergteam}{
    name={ErgTeam},
    sort=ergteam,
    description={Piattaforma di collaborazione aziendale sviluppata da \textit{Ergon Informatica Srl}.}
}

\newglossaryentry{dotnetmauig}{
    name={.NET MAUI},
    sort=dotnetmaui,
    description={Framework di sviluppo multipiattaforma di Microsoft, utilizzato per creare applicazioni native per iOS, Android, Windows e macOS con un unico codice base.}
}

\newglossaryentry{csharpg}{
    name={C\#},
    sort=csharp,
    description={Linguaggio di programmazione multi-paradigma sviluppato da Microsoft, che supporta tutti i concetti della programmazione orientata agli oggetti.}
}

\newglossaryentry{xamlg}{
    name={XAML},
    sort=xaml,
    description={Linguaggio di markup dichiarativo basato su XML sviluppato da Microsoft, utilizzato per definire l'interfaccia utente in modo strutturato, simile a HTML.}
}

\newglossaryentry{apirestg}{
    name={API REST},
    sort=apirest,
    description={Insieme di principi architetturali che permettono la comunicazione tra sistemi client e server attraverso il protocollo HTTP. Esse consentono l'accesso e la manipolazione di risorse mediante l'utilizzo di metodi standard come GET, POST, PUT e DELETE.}
}

\newglossaryentry{mvvmg}{
    name={MVVM},
    text={Model-View-ViewModel},
    sort=mvvm,
    description={Model-View-ViewModel è un pattern architetturale che separa la logica di presentazione dalla logica di business, facilitando il testing e la manutenzione dell'applicazione.}
}

\newglossaryentry{iag}{
    name={IA},
    text={Intelligenza Artificiale},
    sort=intelligenza artificiale,
    description={Intelligenza artificiale è un ramo dell'informatica che si occupa dello sviluppo di sistemi in grado di eseguire compiti che normalmente richiederebbero l'intelligenza umana, come il riconoscimento vocale, la visione artificiale, l'apprendimento automatico e la pianificazione.}
}

\newglossaryentry{umlg}{
    name={UML},
    text={Unified Modeling Language},
    sort=uml,
    description={Unified Modeling Language è un linguaggio di modellazione standardizzato utilizzato per la progettazione di sistemi software.}
}